% !TeX program = xelatex
% !TeX encoding = UTF-8
% !BIB program = bibtex
\documentclass[lettersize,journal]{IEEEtran}
\usepackage{iftex}
\ifXeTeX\else
  \PackageError{ITM}{Please compile main-cn.tex with XeLaTeX}{Select XeLaTeX in TeXstudio or Overleaf.}
\fi
% Fandol ships with TeX Live; no Windows-only fonts are required.
\usepackage[UTF8,fontset=fandol,scheme=plain]{ctex}
\usepackage{amsmath,amssymb}
\usepackage{graphicx}
\usepackage{booktabs,tabularx,array}
\usepackage{seqsplit}
\usepackage{cite}
\usepackage[hidelinks]{hyperref}
\newcolumntype{L}{>{\raggedright\arraybackslash}X}
\setlength{\emergencystretch}{1em}
% \renewcommand{\abstractname}{摘要}
% \renewcommand{\IEEEkeywordsname}{关键词}
% \renewcommand{\refname}{参考文献}
% \renewcommand{\figurename}{图}
% \renewcommand{\tablename}{表}

\begin{document}
\title{ITM：面向文本到动作生成的稀疏 IMU 控制}
\author{作者信息待填写}
\maketitle
\begin{abstract}
% !TeX root = ../main-cn.tex
% UTF-8. Abstract synchronized with paper/论文草稿.md.

文本到人体动作生成能够根据自然语言产生多样动作，但文本难以指定具体运动实例；稀疏惯性测量单元（IMU）提供连续的局部运动线索，却不足以唯一确定全身动作。
本文研究 IMU-guided Text-to-Motion Generation，以文本表达动作意图，以少量 IMU 读数约束生成动作的运动细节。
我们提出面向预训练文本到动作模型的轻量条件适配框架 ITM，将稀疏 IMU 表示、时序控制编码和生成基座适配组织为统一流程。
框架冻结生成骨干与文本编码器，通过可训练的控制模块注入 IMU 信息，以支持向不同生成架构迁移，并以保持文本语义和动作自然性为目标实现实例级运动控制。
我们构建 HumanML3D/AMASS 文本、动作与虚拟 IMU 配对数据，并通过扩散与流匹配基座上的实验考察框架的适用性。
受控条件对照检验 IMU 是否使生成动作更贴近目标运动，文本编辑检验语言能否补充未观测部位的语义，生成质量评价则衡量新增控制对语义匹配和动作分布的影响。
实验揭示了局部运动约束与全身语义生成的互补潜力，以及控制精度、动作平滑性和语义保持之间的权衡。
本研究为将可穿戴传感器接入预训练动作生成器提供了可扩展的适配思路与评价依据。
% 作者注：正文采用面向多基座迁移的设计目标；“支持迁移”不等于任意基座无需适配，更不代表已在最新最强模型上验证。
% 当前已验证范围：MDM、MotionLab；HY-Motion 仅完成 Text-only smoke 和旋转导出，尚无融合训练与评价结果。
% 严谨结果备忘：MDM 多种子结果支持双腕实例级控制，单头收益不显著；不得将该结论推广为所有配置和基座均可靠。
% 严谨结果备忘：摆臂和转向有响应，速度与步幅响应不可靠；控制收益伴随 FID 代价与部分文本补全能力下降。
% 待补：HY-Motion 融合训练、paired/shuffled/zero/Text-only 对照，以及同协议生成质量和多种子控制评价。
% 待补：统一测试 ID、caption、物理时长和重复次数；不同基座的噪声空间不同，仅在同一基座内共享初始噪声。
% 待补：文本编辑同时比较 Text-only 与 Text+IMU 的语义响应和局部约束，检验互补性而非仅检验输出变化。
% 待补：报告各基座适配接口、训练参数和迁移成本；只有具备相应证据后，才将“设计目标”升级为“跨基座稳定有效”的实证结论。

\end{abstract}
\begin{IEEEkeywords}
人体动作生成，文本到动作，稀疏 IMU，可控生成，扩散模型。
\end{IEEEkeywords}

% !TeX root = ../main-cn.tex
% UTF-8. Introduction synchronized with paper/论文草稿.md.
\section{引言}
\label{sec:1_introduction}

人体动作生成连接自然语言理解、人体运动建模和虚拟角色控制。
近年来，文本到动作（Text-to-Motion）模型能够根据自然语言生成多样的人体运动，为动画创作、虚拟角色和沉浸式交互提供了便捷的内容生成方式。
从扩散模型 MDM 到采用流匹配的 MotionLab，不同生成架构持续改善文本语义与动作分布的建模能力。
然而，更强的文本生成能力并不直接意味着用户能够控制某一次动作的具体执行方式。

这一困难源于文本描述的欠定性。
例如，“a person walks”指定了行走这一动作类别，却没有唯一确定步速、步频、转向和身体摆动。
即使补充更详细的描述，用户也难以用语言逐帧表达自身运动的连续变化。
因此，如何在保留文本表达高层意图的便利性的同时，引入可直接反映运动实例的控制信号，是文本动作生成走向交互式应用需要解决的问题。

稀疏惯性测量单元（IMU）为这一问题提供了可穿戴的运动信息来源。
头部或手腕设备的加速度及姿态读数能够反映局部身体运动，为生成过程提供文本之外的连续时序约束。
已有 DIP、TransPose、IMUPoser 和 MobilePoser 等方法主要从少量 IMU 估计真实人体姿态，以重建准确度和运动稳定性为主要评价目标。
但局部观测并不能唯一确定全身动作：相近的头部 IMU 读数可能同时对应自然垂臂行走和明显摆臂行走。
文本可以进一步指定未观测身体部位的动作意图，而生成模型提供满足这些条件的全身运动先验。

本文研究 IMU-guided Text-to-Motion Generation，即在文本与稀疏 IMU 的共同条件下生成完整人体动作。
我们关注两种互补关系：当文本描述模糊时，IMU 能否约束具体运动实例；当身体观测稀疏时，文本能否补充未观测部位的动作语义。
该任务保留随机生成能力，允许同一组条件对应多个合理动作，因此既需要衡量与控制目标的接近程度，也需要评价文本语义、生成多样性和动作自然性。
当文本与传感器信号存在冲突时，两种条件如何共同影响输出，同样是需要分析的任务边界。

实现这种互补并非简单地把两个输入拼接起来。
首先，IMU 描述局部的加速度与朝向，而生成模型通常在关节、旋转或潜在动作表示中建模，两者在坐标系、时间分辨率和物理含义上存在差异。
其次，预训练生成模型已具有较强的文本先验，新增控制可能被忽略，也可能过度改变原有生成分布。
最后，约束观测部位与保留未观测部位的语义自由度可能相互竞争，降低局部误差并不必然改善整体动作质量。
这些问题使得跨基座验证、正确与错配条件的对照，以及多随机种子评价成为检验方法有效性的必要环节。

为此，我们提出面向预训练文本到动作模型的轻量条件适配框架 ITM，将传感器信息处理与生成基座的内部计算分离。
框架由三个相互衔接的环节组成：用统一的稀疏 IMU 表示描述可用传感器及其时序读数，通过控制编码器提取运动特征，再由基座适配层将特征接入生成过程。
不同基座共用条件表示与控制学习思路，适配层负责匹配其特征维度、时间分辨率和条件注入位置。
这种模块划分旨在将迁移工作集中于基座适配环节，使数据构建、条件对照和评价流程能够复用。
训练期间冻结预训练生成骨干与文本编码器，仅优化新增的控制模块，以利用已有动作先验并降低参数更新规模。
以 MDM 为例，基座适配层可通过逐层零初始化残差实现控制注入；这一实现体现了统一框架在特定架构上的实例化方式。
我们同时构建 HumanML3D/AMASS 文本、动作与虚拟 IMU 配对数据，为框架训练和跨基座评价提供基础。

框架的目标是在不同架构与不同生成能力的预训练模型上，以可复用的适配流程获得实例级控制，并保留文本语义与动作自然性。
我们以扩散与流匹配生成器为实验载体，从控制有效性、语义互补性和生成质量三个方面考察这一目标。
正确配对、错配与条件关闭的对照用于检验传感器是否提供来自具体运动实例的信息。
固定 IMU 的文本编辑用于考察语言对未观测身体部位的影响，并与固定文本的 IMU 编辑共同分析两类条件的作用。
生成质量评价、损失消融和重复采样进一步用于区分控制收益、随机变化及对原有生成分布的扰动。
这些实验将基座的文本生成能力与新增 IMU 控制能力分别评价，为框架的扩展和适用范围提供依据。

% 严谨实现说明：统一的是条件表示、冻结基座的学习原则和评价流程；当前不是一个控制器 checkpoint 可直接跨所有基座使用。
% MDM 使用逐层残差适配器；MotionLab 使用原生 hint/control token 接口，编码器维度和实现也有差别。
% 不能声称所有基座都采用逐层零初始化注入，或迁移只需修改一个投影层；这些需通过代码接口与迁移成本验证。
% 已完成的跨基座实验为 MDM 与 MotionLab。HY-Motion 尚无 Text+IMU 结果；待完成训练和统一复测后扩展实证主张。
% 严谨结果备忘：MDM 双腕多种子改善稳定，单头改善不显著；MotionLab 改善率尚未达到预设门槛。
% 严谨结果备忘：摆臂、转向响应较明确，速度与步幅不可靠；更强 Text-only 不保证更强实例控制。
% 待补：相同测试 ID、物理时长和重复协议的近期强基座评价，报告各自 Text-only 与 Text+IMU 的质量及控制差值。
% 待补：统一适配接口的实现清单、可训练参数量、显存与迁移改动范围，以支撑“可复用、可扩展”的工程主张。
% 待补：文本编辑下联合比较语义响应与 active-sensor 误差，加入 Text-only 对照；保留失败案例。

本文的贡献概括如下：

\begin{enumerate}
\item 形式化文本与稀疏 IMU 联合条件下的人体动作生成任务，明确实例级运动控制与未观测部位语义补全这两类互补目标。
\item 提出由统一 IMU 表示、时序控制编码和基座适配组成的轻量框架，将预训练生成能力与传感器控制学习分离，并构建配对虚拟 IMU 数据桥。
\item 建立包含正确配对、错配及条件关闭对照的评价方式，在同一基座内固定初始生成噪声，并结合文本编辑与重复采样区分条件响应和有效控制。
\item 通过不同生成架构上的实验，分析新增传感器控制对语义匹配、实例级约束与动作稳定性的影响，为跨基座迁移提供评价依据。
\end{enumerate}

本文面向“语言指定意图、可穿戴设备提供运动细节”的动作生成场景。
通过把预训练生成能力与传感器控制适配分开设计，我们希望明确稀疏 IMU 可以提供哪些增量信息，以及这些信息在不同生成基座中受到怎样的限制。

% !TeX root = ../main-cn.tex
% UTF-8. Related works synchronized with paper/论文草稿.md.
\section{相关工作}
\label{sec:2_relatedworks}

\subsection{文本到动作生成与预训练基座}

HumanML3D \cite{guo2022humanml3d} 为文本到人体动作生成提供了配对数据与常用评价协议，推动了连续动作建模、潜在空间生成和离散动作 token 建模等路线的发展。
MDM \cite{tevet2023mdm} 与 MotionDiffuse \cite{zhang2024motiondiffuse} 使用扩散过程建模文本条件下的动作分布；MLD \cite{chen2023mld} 则先学习低维动作表示，再在潜在空间中进行扩散，以降低生成的计算开销。
T2M-GPT \cite{zhang2023t2mgpt}、MotionGPT \cite{jiang2023motiongpt} 和 MoMask \cite{guo2024momask} 从离散表示出发，分别探索自回归、语言与动作联合建模以及掩码 token 生成。
LGTM \cite{sun2024lgtm} 研究局部到全局的文本驱动动作建模，进一步增强了细粒度身体语义与整体动作结构之间的联系。

流匹配也为动作生成提供了另一类生成建模范式。
MotionLab \cite{guo2025motionlab} 通过 Motion-Condition-Motion 范式和 rectified flow 统一动作生成与编辑，而 HY-Motion 1.0 \cite{tencent2025hymotion} 探索了大规模流匹配模型在文本动作生成中的扩展能力。
总体而言，现有文本动作生成方法已经形成了扩散、潜在扩散、离散 token 和流匹配等多种成熟基座，为进一步引入外部控制信号提供了较强的动作先验。本文关注如何在尽量保留这些预训练先验的同时，将稀疏 IMU 观测作为额外条件接入生成过程。

\subsection{可控动作生成与条件适配}

可控动作生成在文本语义之外加入轨迹、关键帧或局部关节约束，使用户能够指定动作的空间执行方式。
GMD \cite{karunratanakul2023gmd} 通过特征投影与引导机制结合文本和空间约束，处理轨迹及稀疏关键帧等控制条件。
OmniControl \cite{xie2023omnicontrol} 进一步支持不同时间与不同关节的空间约束，并结合空间引导和真实感引导协调控制精度与动作自然性。
这些工作表明，额外条件能够显著提高动作生成的可控性，同时也需要兼顾条件满足程度与生成动作的自然性。

在预训练模型适配方面，ControlNet \cite{zhang2023controlnet} 在图像扩散模型上展示了通过零初始化控制连接引入额外条件、同时保留原有生成先验的有效性，这一思路也为动作生成中的条件适配提供了参考。
近期 MoCoDiff \cite{song2026mocodiff} 使用模态相关的轻量调制模块处理文本、风格和历史动作条件，进一步关注多条件之间的干扰以及长时序一致性。
因此，将轻量条件模块接入预训练生成模型已经成为可控生成中的常见思路，而不同控制信号的物理含义和信息密度会进一步影响条件表示与融合方式。

与直接给定关节位置或轨迹不同，IMU 通常提供局部加速度和朝向信息，并不显式包含完整的绝对人体轨迹；同时，传感器安装坐标、初始状态和观测部位也会影响可恢复的信息。
这使得稀疏 IMU 控制既包含连续时序约束，又具有明显的局部性和不完备性。
本文据此采用时序表示将 IMU 信号接入文本动作生成过程，并通过不同条件组合考察文本语义与实例级惯性观测之间的互补关系。

\subsection{稀疏惯性人体动作重建}

稀疏惯性动作重建利用少量可穿戴设备估计全身姿态，为低成本人体运动捕捉提供了技术基础。
DIP \cite{huang2018dip} 学习从稀疏惯性观测恢复人体姿态，TransPose \cite{yi2021transpose} 进一步联合估计姿态与人体平移。
PIP \cite{yi2022pip} 在稀疏惯性估计中进一步引入物理约束，将神经运动学估计与物理优化结合，以改善长时序运动的稳定性和物理合理性。
IMUPoser \cite{mollyn2023imuposer} 和 MobilePoser \cite{xu2024mobileposer} 则将这一方向拓展到手机、手表和耳戴设备等消费级硬件，强调稀疏设备条件下的实际可用性。
这些工作为传感器表示、时序建模、全局运动估计以及人体物理约束提供了重要基础。

稀疏惯性重建与本文研究在输入模态上高度相关，但二者的任务目标并不完全相同。
传统重建方法通常以对应真实运动为唯一目标，重点考察姿态、平移和时间稳定性；文本条件动作生成则同时涉及语义匹配、动作分布以及同一语义下的合理多样性。
因此，惯性重建方法能够反映模型从稀疏传感器中恢复人体运动信息的能力，而文本与 IMU 联合生成还需要进一步考察两类条件共同作用时的生成质量和条件一致性。

\subsection{语言与传感器的联合动作建模}

语言与稀疏传感器联合建模与本文最为接近。
Spatial-Related Sensors Matters \cite{yang2024spatialsensors} 利用文本语义辅助惯性动作重建，通过传感器加权、层次时序 Transformer 和对比学习处理传感器特征与语义的对齐。
该工作表明，文本语义能够帮助缓解稀疏惯性观测带来的动作歧义，体现了语言与惯性信号之间的互补性。

Ego4o \cite{wang2025ego4o} 接收不同组合的稀疏 IMU、可选第一视角图像与动作描述，将多模态信息映射到动作 VQ-VAE 的潜在空间，用于动作捕捉与理解。
在不使用视觉输入时，其文本与 IMU 的联合设置与本文具有较直接的输入输出联系。
EgoLM \cite{hong2025egolm} 则借助语言模型联合建模动作与语言，覆盖动作描述以及从文本或稀疏传感器生成动作等任务。
这些研究已经开始将语言、惯性传感器和人体动作放入统一建模框架中，说明多模态条件不仅可以辅助运动重建，也可以用于更一般的动作生成与理解任务。

与面向多模态、多任务的统一系统相比，本文更聚焦于文本与稀疏 IMU 两类条件在预训练文本动作生成基座中的协同作用。
文本主要描述动作类别和语义意图，而 IMU 提供与具体执行实例相关的局部动态观测；二者分别对应较高层的语义约束和较低层的时序运动约束。
因此，如何在保留文本动作先验的同时利用稀疏惯性观测，是本文所研究问题与已有语言辅助惯性重建和多模态动作建模工作的主要联系所在。

% !TeX root = ../main-cn.tex
% Converted from paper/论文草稿.md; UTF-8.
\section{方法}
\label{sec:3_method}


\subsection{Problem Formulation}

给定文本描述 $c$、稀疏 IMU 序列 $s$、传感器有效 mask $m$ 和随机噪声 $z$，ITM 学习条件分布 $p_\theta(x\mid c,s,m)$，并生成动作样本 $x=G_\theta(c,s,m,z)$。
文本描述动作语义，IMU 提供局部身体部位的连续时序控制，mask 指明实际可用的传感器。
由于随机噪声允许同一组条件产生多个合理动作，任务目标不是确定性恢复唯一 ground truth，而是在保持文本语义和动作自然性的同时响应给定 IMU。

动作采用 HumanML3D 263D representation。
该表示联合编码根节点运动、局部关节位置、关节旋转、关节速度和足部接触；生成后使用 HumanML3D 的 \texttt{\seqsplit{recover\_from\_ric}} 过程积分根节点运动，并把局部关节坐标转换为 22 个全局三维关节。
IMU表示为6个固定槽位，每个槽位每帧包含3D acceleration和3x3 orientation，共12D。
方向采用完整旋转矩阵，是为了直接保存传感器坐标系朝向并与Slerp重采样兼容；将其改为连续6D旋转表示是可行的输入消融，但不属于当前实现。

\subsection{Standard Virtual IMU Bridge}

我们通过 HumanML3D \texttt{\seqsplit{index.csv}} 将 HumanML3D motion 映射回 AMASS/SMPL 序列，并从 SMPL/AMASS 生成标准虚拟 IMU。
缓存频率为 30 FPS，训练和采样时重采样到 MDM 的 20 FPS。
Acceleration 使用线性插值，orientation 使用 Slerp。
每个 IMU 序列保留 6 个固定槽位：

\begin{table}[t]
\centering
\caption{虚拟 IMU 的固定传感器槽位}
\label{tab:result-1}
\footnotesize
\setlength{\tabcolsep}{3pt}
\renewcommand{\arraystretch}{1.15}
\begin{tabularx}{\linewidth}{LL}
\toprule
Slot & Part \\
\midrule
0 & Left Wrist \\
1 & Right Wrist \\
2 & Left Thigh \\
3 & Right Thigh \\
4 & Head \\
5 & Pelvis \\
\bottomrule
\end{tabularx}
\end{table}


融合模型训练使用两种配置：head-only，即slot 4；wrists，即slots 0和1。
头部配置对应头显、耳机和智能眼镜等常见设备，也是传感器最少、全身动作最欠定的设置；双腕配置对应智能手表或手环，能够提供更直接的上肢控制信号。
两者分别用于检验“文本补全未观测身体”和“IMU控制局部运动”两类互补现象。
数据规模为7009 train、434 val、1333 test aligned records。

\subsection{Frozen MDM Backbone}

我们使用官方MDM HumanML3D checkpoint \texttt{\seqsplit{humanml\_trans\_enc\_512/model000475000.pt}}。
该模型使用CLIP ViT-B/32编码文本，Transformer latent width为512，包含8层Transformer Encoder，diffusion steps为1000，最大序列长度为196 frames。
训练期间冻结MDM和CLIP，不更新其参数。

冻结 MDM 的目的有三点。
第一，保留官方模型已经学到的文本到动作生成先验。
第二，降低训练成本，使实验可在单卡上快速迭代。
第三，保证初始化时 Text-only 路径与官方 MDM 一致，便于分析 IMU 控制是否真正产生增量作用。

\subsection{Sparse IMU Encoder}

IMU encoder 输入为：

\[s\in\mathbb{R}^{B\times T\times6\times12}.\]
$B$表示batch size，$T$表示动作帧数，6为固定传感器槽位数，12为每个传感器的3维加速度和9维方向矩阵。
每个 12D sensor feature 先通过共享线性层投影到 512D，并加入可学习 sensor embedding。
随后使用 sensor mask 在传感器维度做 masked mean pooling，得到每帧一个 512D token。
该序列再输入 2 层、8 heads 的 temporal Transformer Encoder，得到 IMU control feature：

\[C_{\mathrm{imu}}\in\mathbb{R}^{B\times T\times512}.\]
$C_{imu}$是逐帧控制序列，其中每个时间步由一个512维向量概括当前有效传感器及其前后时序上下文。
本文在时序编码前进行多传感器pooling。
这使模型能够直接支持可变sensor mask，但也可能削弱不同传感器之间的显式交互。

\subsection{Zero-initialized Residual Adapters}

我们在 MDM 的 8 个 Transformer layer 后分别添加一个零初始化线性 adapter：

\begin{align}\widetilde H_l &= \operatorname{MDMLayer}_l(H_{l-1}),\\ H_l &= \widetilde H_l + W_l C_{\mathrm{imu}}+b_l.\end{align}

每层 adapter 独立，不共享参数。
由于 adapter 权重和 bias 都初始化为 0，初始模型与原始 MDM 完全等价。
训练过程中，IMU encoder 学习提取控制特征，adapter 学习将这些特征注入不同层的 motion tokens。

需要注意的是，MDM 内部序列包含一个 text/timestep condition token 和 T 个 motion tokens，而 IMU 只对应 T 个时间帧。
因此我们在 control sequence 前补一个 zero token，使其与 MDM 的 T+1 token 对齐。

\subsection{Training Objective}

基础融合模型（工程记录名Stage-1）使用官方 MDM diffusion loss，即 predicted x\_start reconstruction loss。
训练时文本条件和 IMU 条件独立 dropout：

\begin{itemize}
\item text dropout probability: 0.1。
\item IMU dropout probability: 0.1。
\end{itemize}

训练超参数如下：

\begin{table}[t]
\centering
\caption{训练配置}
\label{tab:result-2}
\footnotesize
\setlength{\tabcolsep}{3pt}
\renewcommand{\arraystretch}{1.15}
\begin{tabularx}{\linewidth}{LL}
\toprule
Item & Value \\
\midrule
Training records & 7009 \\
Sensor configs & head, wrists \\
Epochs & 5 \\
Batch size & 16 \\
Optimizer & AdamW \\
Learning rate & 1e-4 \\
Gradient clipping & 1.0 \\
Frozen modules & MDM, CLIP \\
Trainable modules & IMU encoder, 8 adapters \\
\bottomrule
\end{tabularx}
\end{table}


五个 epoch 的 loss 为：

\begin{quote}0.07791, 0.08080, 0.07893, 0.07619, 0.07642\end{quote}

基础融合模型不包含显式 IMU consistency、velocity consistency、joint trajectory、foot contact 或 physics loss。
实验中该模型能够响应 IMU，但 jerk 偏高，因此我们进一步引入轻量的关节空间正则。

在正则化主模型中，我们保持MDM和CLIP冻结，从基础融合checkpoint继续训练IMU encoder和adapters。
新增辅助损失包括active sensor trajectory loss、active sensor velocity loss和jerk regularization。
本文主模型（实验记录名\texttt{\seqsplit{stage2b\_balanced\_b}}）使用如下设置：

\begin{table}[t]
\centering
\caption{训练配置}
\label{tab:result-3}
\footnotesize
\setlength{\tabcolsep}{3pt}
\renewcommand{\arraystretch}{1.15}
\begin{tabularx}{\linewidth}{LL}
\toprule
Item & Value \\
\midrule
Resume checkpoint & \texttt{\seqsplit{stage1\_full\_pilot\_v2.pt}} \\
Output checkpoint & \texttt{\seqsplit{stage2b\_balanced\_b.pt}} \\
Epochs & continue to 8 \\
Batch size & 16 \\
Learning rate & 5e-5 \\
Trajectory loss weight & 1.0 \\
Velocity loss weight & 0.1 \\
Jerk loss weight & 0.003 \\
Frozen modules & MDM, CLIP \\
Trainable modules & IMU encoder, 8 adapters \\
\bottomrule
\end{tabularx}
\end{table}


辅助损失先将 normalized HumanML3D 263D motion 反标准化，再通过 torch 版 \texttt{\seqsplit{recover\_from\_ric}} 恢复 22 joints。
Active sensor joints 按配置选择：head 对应 joint 15，wrists 对应 joints 20 和 21。
该设计不引入 SMPL decoder、cross-attention 或完整 ControlNet backbone copy。

作为方法消融，我们还加入 frozen MDM 的 Text-only prediction作为anchor，希望缓解正则化后“更平滑但文本补全变弱”的trade-off。
普通text anchor loss权重为0.2，upper-body text anchor loss权重为0.4；后者只约束非active-sensor的上肢关节。
该消融略微恢复same-head arm swing和active head error，但仍未达到基础融合模型的文本补全强度，因此不作为主模型。
工程文档中的Stage-1、Stage-2b和Stage-4分别对应基础融合、正则化主模型和text-anchor消融；正文后续仅在消融表中保留这些简称。

\subsection{Independent Guidance}

推理时我们使用三分支 guidance：

\begin{equation}\begin{split}f_{\mathrm{guided}}={}&f_{00}+w_{\mathrm{text}}(f_{10}-f_{00})\\ &+w_{\mathrm{imu}}(f_{11}-f_{10}).\end{split}\end{equation}
式中，$f_{00}$、$f_{10}$ 和 $f_{11}$ 分别表示无条件、仅文本和文本与 IMU 联合条件下的预测；$w_{\mathrm{text}}$ 和 $w_{\mathrm{imu}}$ 分别对应 text scale 和 IMU scale。
其中\texttt{\seqsplit{unconditional}}同时关闭文本和IMU，\texttt{\seqsplit{text\_only}}只启用文本，\texttt{\seqsplit{text\_imu}}同时启用文本和IMU。
第一项差分表示文本相对无条件生成带来的变化，第二项差分表示IMU在给定文本后的增量控制。
这使得 Text-only、IMU-only diagnostic、Text+IMU 可以通过同一个模型和同一组噪声比较。
受控条件交换实验中所有条件共享 diffusion initial noise，以确保可见差异来自条件变化，而不是随机噪声变化。

% !TeX root = ../main-cn.tex
% Converted from paper/论文草稿.md; UTF-8.
\section{实验}
\label{sec:4_experiments}


\subsection{Datasets}

主实验使用 HumanML3D/AMASS 对齐数据。
HumanML3D 提供文本和 263D motion；AMASS/SMPL 提供生成标准虚拟 IMU 所需的身体运动。
标准 IMU cache 包含：

\begin{table}[t]
\centering
\caption{对齐数据集划分}
\label{tab:result-4}
\footnotesize
\setlength{\tabcolsep}{3pt}
\renewcommand{\arraystretch}{1.15}
\begin{tabularx}{\linewidth}{LL}
\toprule
Split & Records \\
\midrule
Train & 7009 \\
Val & 434 \\
Test & 1333 \\
\bottomrule
\end{tabularx}
\end{table}


早期 sanity baseline 使用 999/200/200 的 train/val/test 子集，用于验证数据流、回归 baseline 和 sparse sensor ablation。

我们还在Nymeria上进行小规模真实IMU预备实验。
Nymeria的\texttt{\seqsplit{body/xdata.npz}}包含Xsens全身动作与17个真实IMU，\texttt{\seqsplit{motion\_narration.csv}}提供分段动作描述。
该实验仅使用3条来自不同录制序列、时长约5秒且head/wrists四元数有效的片段，不参与HumanML3D主表，也不用于训练或选择主模型。

\subsection{Metrics}

本文使用三类指标。

第一类是Text-to-Motion标准指标。
IMU-mapped subset表报告Matching Score、R-precision、FID和Diversity；Multimodality列入定义，但未在该子集协议中报告。

\begin{table}[t]
\centering
\caption{评价指标及其含义}
\label{tab:result-5}
\footnotesize
\setlength{\tabcolsep}{3pt}
\renewcommand{\arraystretch}{1.15}
\begin{tabularx}{\linewidth}{LL}
\toprule
Metric & Meaning \\
\midrule
Matching Score & 文本与动作 embedding 距离，越低越好 \\
R-precision & 文本-动作检索 top-k 准确率，越高越好 \\
FID & 生成动作分布与真实分布距离，越低越好 \\
Diversity & 生成动作整体多样性 \\
Multimodality & 同一文本下多样本差异 \\
\bottomrule
\end{tabularx}
\end{table}


第二类是受控条件交换实验使用的控制与动作质量指标：

\begin{table}[t]
\centering
\caption{评价指标及其含义}
\label{tab:result-6}
\footnotesize
\setlength{\tabcolsep}{3pt}
\renewcommand{\arraystretch}{1.15}
\begin{tabularx}{\linewidth}{LL}
\toprule
Metric & Meaning \\
\midrule
active\_sensor\_trajectory\_error\_m & 有效传感器对应joint的root-relative trajectory error \\
active\_sensor\_acceleration\_error\_mps2 & 由生成active joints二阶差分得到的加速度与目标IMU加速度之间的proxy误差 \\
head\_trajectory\_error\_m & head joint 轨迹误差 \\
wrist\_trajectory\_error\_m & wrists 轨迹误差 \\
root\_relative\_motion\_error\_m & 全身 root-relative joint error \\
jerk\_ratio & 生成 jerk / GT jerk \\
arm\_swing\_proxy\_m & wrist-relative-to-shoulder motion amplitude \\
root\_travel\_m & root 水平位移 \\
step\_frequency\_hz & foot speed frequency proxy \\
\bottomrule
\end{tabularx}
\end{table}


第三类是 IMU-only baseline 指标，包括 root-relative MPJPE、rotation geodesic error、jerk ratio 和 foot skating。
这类指标适合 IMU pose estimation，不应作为 ITM 的唯一主指标。

Active-sensor trajectory error首先将预测与GT转换为root-relative joint trajectories，再计算有效传感器对应关节的平均欧氏距离。
Acceleration error由HumanML 22-joint轨迹二阶差分得到；由于当前尚未将生成结果拟合为SMPL并恢复传感器安装坐标，该值不是严格的generated-IMU acceleration误差，更不包含orientation geodesic error。
正式传感器一致性评价应使用数据桥中已有的虚拟传感器定义，将生成动作恢复到SMPL后重新导出传感器加速度和方向；本文现阶段只把关节空间量用于同协议模型诊断。

对比实验采用分层协议。
第一层与官方MDM Text-only比较生成质量，回答加入IMU后文本生成能力是否退化。
第二层以Flexible IMUPoser描述IMU-only reconstruction boundary；它没有文本输入和随机生成目标，因此不进入FID/R-precision表。
第三层比较Stage-1、Stage-2b和Stage-4，分析平滑性、IMU控制与head-only文本补全的trade-off。
Ego4o和Spatial-Related Sensors Matters与本文最接近，但缺少可直接核验的官方checkpoint或实现，本文不填入无法复现的数字。

\subsection{Diagnostic Reconstruction Baselines}

早期线性 baseline 表明，synthetic IMU-heavy 输入在 deterministic reconstruction 中很强，而文本特征帮助有限。
999 train / 200 test 上，IMU-only linear baseline 的 test MSE 为 0.024277，text-only 为 0.062518。

Frame-level MLP 在 999 train / 200 val 上进一步降低 IMU-only eval MSE 到 0.014230，但加入 hashed text 后 eval MSE 反而恶化到 0.020078。
这说明简单文本表示容易记忆训练集，不等于文本语义真正有用。

Temporal Transformer + DistilBERT 将文本换成预训练 embedding，并加入序列建模，但在 6 sensors 条件下，IMU-only test MSE 为 0.018271，DistilBERT+IMU 为 0.018774，文本仍未改善。
这提示：在强 IMU 和 deterministic MSE 目标下，文本难以发挥作用。

\subsection{Sparse Sensor Ablation of the Diagnostic Model}

稀疏传感器实验显示，文本价值依赖传感器位置和欠定程度。
在 wrists 配置下，DistilBERT+IMU 相比 IMU-only 同时改善 test MSE 和 MAE：

\begin{table}[t]
\centering
\caption{Sparse Sensor Ablation of the Diagnostic Model：实验结果}
\label{tab:result-7}
\footnotesize
\setlength{\tabcolsep}{3pt}
\renewcommand{\arraystretch}{1.15}
\begin{tabularx}{\linewidth}{LLLL}
\toprule
Sensors & Text & Test MSE & Test MAE \\
\midrule
wrists & no & 0.030797 & 0.098079 \\
wrists & yes & 0.029610 & 0.096633 \\
\bottomrule
\end{tabularx}
\end{table}


但 head-only 的验证集改善没有在 test 上复现；6 sensors 仍 favor IMU-only。
这说明旧回归任务并不适合作为最终 formulation，也促使我们转向生成式 MDM backbone。

\subsection{Text-only MDM Baseline}

我们接入官方 MDM checkpoint 作为合格 Text-only baseline。
官方 20-rep reference log 报告：

\begin{table}[t]
\centering
\caption{评价指标及其含义}
\label{tab:result-8}
\footnotesize
\setlength{\tabcolsep}{3pt}
\renewcommand{\arraystretch}{1.15}
\begin{tabularx}{\linewidth}{LL}
\toprule
Metric & Value \\
\midrule
FID & 0.5443 \\
R-precision top-1 & 0.3195 \\
R-precision top-2 & 0.4978 \\
R-precision top-3 & 0.6110 \\
Diversity & 9.5595 \\
\bottomrule
\end{tabularx}
\end{table}


ITM evaluator bridge复用官方HumanML3D text-motion evaluator wrapper，对生成结果计算Matching Score、R-precision、FID和Diversity。
IMU-mapped test subset共有682条可用动作；为保持batch size 32，实际评价前672条，并独立重复5次。
该表用于同一子集上的直接比较，不与full-set 20-rep文献数字混排。

结果如下：

\begin{table*}[t]
\centering
\caption{Text-only MDM Baseline：实验结果}
\label{tab:result-9}
\footnotesize
\setlength{\tabcolsep}{3pt}
\renewcommand{\arraystretch}{1.15}
\begin{tabularx}{\linewidth}{LLLLLLLL}
\toprule
Model & Sensor & Matching \ensuremath{\downarrow} & R@1 \ensuremath{\uparrow} & R@2 \ensuremath{\uparrow} & R@3 \ensuremath{\uparrow} & FID \ensuremath{\downarrow} & Diversity \\
\midrule
MDM Text-only subset & none & 3.4152 \ensuremath{\pm} 0.0196 & 0.4530 \ensuremath{\pm} 0.0178 & 0.6506 \ensuremath{\pm} 0.0157 & 0.7586 \ensuremath{\pm} 0.0116 & 0.5089 \ensuremath{\pm} 0.0441 & 10.0517 \ensuremath{\pm} 0.1657 \\
ITM Stage-2b & head & 3.2680 \ensuremath{\pm} 0.0519 & 0.4580 \ensuremath{\pm} 0.0131 & 0.6682 \ensuremath{\pm} 0.0187 & 0.7780 \ensuremath{\pm} 0.0139 & 0.6526 \ensuremath{\pm} 0.0521 & 10.1686 \ensuremath{\pm} 0.1396 \\
ITM Stage-2b & wrists & 3.3546 \ensuremath{\pm} 0.0378 & 0.4551 \ensuremath{\pm} 0.0150 & 0.6598 \ensuremath{\pm} 0.0218 & 0.7673 \ensuremath{\pm} 0.0103 & 0.8167 \ensuremath{\pm} 0.0357 & 10.0319 \ensuremath{\pm} 0.1746 \\
\bottomrule
\end{tabularx}
\end{table*}


Stage-2b head的Matching Score和R-precision与Text-only接近，但FID从0.5089增加到0.6526；wrists的检索指标同样接近Text-only，FID则增加到0.8167。
这表明新增条件没有造成文本语义崩溃，但会扰动MDM原有生成分布，且双腕控制的分布代价更明显。
Stage-4未在该协议下评价，因此仅进入控制消融表。

为检验方法能否迁移到更强的Text-to-Motion骨干，我们进一步复现MotionLab MotionFlow，并冻结其MotionFlow与CLIP，只训练IMU时序编码器。
该迁移使用MotionFlow原生512维hint token流，而不是把IMU伪装成66维关节轨迹。
相同672条子集和5次重复的结果如下：

\begin{table*}[t]
\centering
\caption{Text-only MDM Baseline：实验结果}
\label{tab:result-10}
\footnotesize
\setlength{\tabcolsep}{3pt}
\renewcommand{\arraystretch}{1.15}
\begin{tabularx}{\linewidth}{LLLLLLL}
\toprule
模型 & Matching Score\ensuremath{\downarrow} & R@1\ensuremath{\uparrow} & R@2\ensuremath{\uparrow} & R@3\ensuremath{\uparrow} & FID\ensuremath{\downarrow} & Diversity\ensuremath{\rightarrow} \\
\midrule
MotionLab Text-only & 2.7223\ensuremath{\pm}0.0230 & 0.5292 & 0.7438 & 0.8321 & 0.2683\ensuremath{\pm}0.0217 & 9.8051 \\
MotionLab + head IMU & 2.8010\ensuremath{\pm}0.0348 & 0.5170 & 0.7280 & 0.8193 & 0.3173\ensuremath{\pm}0.0239 & 10.4175 \\
MotionLab + wrists IMU & 2.8025\ensuremath{\pm}0.0364 & 0.5196 & 0.7185 & 0.8205 & 0.3879\ensuremath{\pm}0.0472 & 10.1779 \\
\bottomrule
\end{tabularx}
\end{table*}


MotionLab Text-only明显强于本文使用的MDM底座。
head IMU条件下Matching、R@3和FID相对退化约2.9\%、1.5\%和18.3\%，说明文本生成质量基本保留；wrists的FID增幅约44.6\%，显示更强局部控制会显著扰动生成分布。
由于后述实例级控制改善率仍未达到预设门槛，本文不将MotionLab迁移模型替换为主模型，而将其作为跨backbone可迁移性与任务难度的诊断。

\subsection{IMU-only Baseline}

Flexible IMUPoser baseline 使用 2-layer 512-wide bidirectional LSTM，输出 24 SMPL joint rotations in 6D，并支持 flexible sensor mask。
test split 上：

\begin{table*}[t]
\centering
\caption{IMU-only Baseline：实验结果}
\label{tab:result-11}
\footnotesize
\setlength{\tabcolsep}{3pt}
\renewcommand{\arraystretch}{1.15}
\begin{tabularx}{\linewidth}{LLLLL}
\toprule
Configuration & MPJPE & Rotation error & Jerk ratio & Foot skating \\
\midrule
head & 14.91 cm & 0.278 rad & 0.294 & 0.123 m/s \\
wrists & 10.74 cm & 0.233 rad & 0.455 & 0.308 m/s \\
\bottomrule
\end{tabularx}
\end{table*}


Wrists 已通过初步 12 cm 门槛；head-only 尚未通过。
该 baseline 用于说明 IMU-only reconstruction 的能力边界，不直接进入 ITM 的文本匹配、FID 或多样性排名。
我们报告其重建指标，但不把这些数字包装成与ITM的胜负关系：IMUPoser估计唯一真实姿态，ITM则在文本语义和随机扩散采样下生成受IMU控制的合理动作。
只有在统一输出表示、关闭生成随机性并明确采用重建目标时，MPJPE才适合作为二者的直接比较。

\subsection{Condition Ablation}

独立条件dropout和guidance使同一个ITM模型支持Text-only、IMU-only diagnostic与Text+IMU。
历史实验默认采用三分支legacy guidance，此时Text-only通过令\texttt{\seqsplit{imu\_scale=0}}得到。
为严格分离IMU主效应与Text-IMU交互，我们另实现四分支factorized guidance；在该模式下，严格Text-only必须同时令\texttt{\seqsplit{imu\_scale=0, joint\_scale=0}}，IMU-only则关闭文本主效应并保留IMU分支。
Text+IMU同时启用两种条件。
所有受控条件交换比较共享动作长度、seed与初始扩散噪声，因而条件间差异不来自随机采样变化。

Text-only负责检验融合模型是否保留MDM语义先验；IMU-only diagnostic观察传感器可提供的轨迹与节奏，但不期待其在稀疏配置下恢复完整动作类型；Text+IMU则用于检验两者是否产生增量互补。
早期Stage-1 pilot已观察到条件交换信号，但也暴露出高jerk，因此主实验进一步比较Stage-1、Stage-2b和Stage-4。

为排除不同阶段训练设置造成的混杂，我们还从同一Stage-1 checkpoint出发，以相同数据顺序、seed、训练轮数和legacy guidance训练四个严格受控变体。
表中每个变体均覆盖60个run和770个generated cases：

\begin{table*}[t]
\centering
\caption{严格受控的辅助损失消融}
\label{tab:result-12}
\footnotesize
\setlength{\tabcolsep}{3pt}
\renewcommand{\arraystretch}{1.15}
\begin{tabularx}{\linewidth}{LLLLLLL}
\toprule
Loss & Same-text trajectory & Same-text jerk & Same-head trajectory & Same-head jerk & Same-head arm swing & Matrix trajectory \\
\midrule
diffusion only & \textbf{0.2725 m} & 1.1782 & 0.1975 m & 2.3268 & \textbf{0.1108 m} & \textbf{0.3597 m} \\
+ trajectory & 0.2764 m & \textbf{0.9989} & \textbf{0.1910 m} & \textbf{1.9784} & 0.0940 m & 0.3657 m \\
+ trajectory + velocity & 0.2782 m & 1.0240 & 0.1932 m & 2.0913 & 0.0992 m & 0.3734 m \\
+ trajectory + velocity + jerk & 0.2766 m & 1.0237 & 0.1922 m & 2.0219 & 0.1026 m & 0.3693 m \\
\bottomrule
\end{tabularx}
\end{table*}


Trajectory loss明显降低了same-text与same-head的jerk，并改善same-head tracking，但同时降低arm swing；在same-text和matrix协议中，辅助损失没有改善active trajectory error。
Velocity与jerk项也没有稳定超过trajectory-only。
因此该严格消融支持“平滑性与文本补全存在trade-off”，但不支持“叠加更多辅助损失会单调改善IMU控制”的结论。

\subsection{Controlled Condition-Swapping Suite on Test Split}

我们在test split上完成60个run、770个generated cases，包括同文本换IMU、同head IMU换文本、guidance sweep以及20组Text/IMU A-B矩阵。

\subsubsection{Same Text, Different IMU}

固定文本：

\begin{quote}a person walks\end{quote}

对30个test samples生成MDM Text-only、ITM Text+head IMU、ITM Text+wrist IMUs。
Stage-1、Stage-2b主模型与Stage-4 text-anchor ablation对比如下：

\begin{table*}[t]
\centering
\caption{Controlled Condition-Swapping Suite on Test Split：实验结果}
\label{tab:result-13}
\footnotesize
\setlength{\tabcolsep}{3pt}
\renewcommand{\arraystretch}{1.15}
\begin{tabularx}{\linewidth}{LLLLLLL}
\toprule
Model & Active sensor error & Active acceleration error & Jerk ratio & Arm swing & Root travel & Step frequency \\
\midrule
Stage-1 & 0.2871 m & 4.7340 m/s\textasciicircum{}2 & 4.2072 & 0.1909 m & 2.1561 m & 1.5505 Hz \\
Stage-2b & 0.2833 m & 4.5360 m/s\textasciicircum{}2 & 1.7599 & 0.1615 m & 1.4666 m & 1.4815 Hz \\
Stage-4 text-anchor & 0.2812 m & 4.6077 m/s\textasciicircum{}2 & 1.8050 & 0.1675 m & 1.4415 m & 1.4695 Hz \\
\bottomrule
\end{tabularx}
\end{table*}


结果表明，同一句模糊文本下更换IMU会改变root travel、step frequency和full-body motion proxies。
这支持IMU提供具体运动约束的假设。
Stage-2b在基本保持active sensor error的同时将jerk ratio从4.2072降至1.7599，并将active acceleration proxy从4.7340 m/s\textasciicircum{}2降至4.5360 m/s\textasciicircum{}2，因此被选为IMU-control主模型。
该acceleration proxy由HumanML 22 joints计算，不等价于完整generated-IMU orientation consistency。

为在更广动作分布上检验这一结论，我们进一步让MDM Stage-2b与MotionLab共享固定100条test motions，并比较严格Text-only、paired、zero和shuffled IMU。
Factorized guidance中的严格Text-only同时关闭IMU主效应和Text-IMU交互项，即\texttt{\seqsplit{imu\_scale=0, joint\_scale=0}}；head和wrists两次Text-only输出完全一致。
Head paired相对Text-only的active error降低$0.0110$ m（95\% CI $[-0.0204,-0.0010]$），但相对shuffled仅降低$0.0017$ m且不显著，说明head的实例级控制仍弱。
Wrists paired相对Text-only、zero和shuffled分别降低$0.0196$ m、$0.0101$ m和$0.0203$ m，paired优于shuffled达到74/100（95\% CI 65\%--82\%，$p<0.0001$），首次超过预设70\%稳定门槛。
相比之下，该通用100-motion子集上的jerk差值区间跨零，因此前述30条walking反事实实验中的Stage-2b平滑收益不能直接泛化到所有动作类别。

MotionLab迁移模型也在同一固定100条test样本上比较Text-only、paired、zero和shuffled IMU。
Head paired将平均active error从0.1733 m降至0.1548 m，配对差值为$-0.0185$ m（95\% bootstrap CI $[-0.0319,-0.0073]$，$p=0.0010$）；wrists从0.3924 m降至0.3615 m，差值为$-0.0309$ m（95\% CI $[-0.0512,-0.0122]$，$p=0.0024$）。
paired相对shuffled的误差也显著更低，head与wrists差值分别为$-0.0070$ m和$-0.0159$ m。
与此同时，paired相对zero control的差异均不显著，说明模型不仅利用了实例信息，也仍受到“开启控制通道”固定偏置的影响。
paired优于Text-only的样本仅为60/100和63/100，其改善率95\%区间分别为50\%--69\%和54\%--72\%，均未稳定达到预设70\%门槛。
Jerk ratio从2.319降至2.005和1.971，对应配对差值的95\%区间均严格低于零。
上述结果支持“平均控制和平滑收益”，但不支持“绝大多数样本均可靠跟随IMU”的更强主张。

我们进一步对该共享100-motion subset进行探索性分层。
语义类别使用固定关键词多标签规则，因类别重叠和小样本量而不作显著性主张。
MDM wrists在locomotion中有26/34条优于Text-only、28/34条优于shuffled，在turning中分别为13/14和11/14，是最稳定的两个类别。
MDM head在upper-body类别的平均tracking变化接近零，jerk反而增加0.240。
单seed长度分组曾显示MDM head在短序列上恶化，但该现象需要通过下述多seed实验复核。

为检验结论是否依赖扩散噪声，我们从同一subset按长度分层抽取30条动作（短、中、长各10条），使用5个diffusion seeds重复严格Text-only、paired和固定shuffled条件。
统计时先对每条motion的5个seed求平均，再以motion为bootstrap单位。
Head paired相对Text-only和shuffled的误差变化分别为$-0.0038$ m（95\% CI $[-0.0108,0.0031]$）和$-0.0013$ m（$[-0.0088,0.0056]$），仅13/30和15/30条动作在至少3/5 seeds上改善。
Wrists则分别降低$0.0281$ m（$[-0.0414,-0.0160]$）和$0.0179$ m（$[-0.0296,-0.0078]$），24/30和22/30条动作在多数seed上改善，且五个seed的平均差异方向完全一致。
按长度看，wrists相对Text-only在短、中、长三组的区间均严格小于零；head各长度组tracking区间均跨零。
该实验支持“双腕IMU提供跨随机种子的实例级控制”，但不支持对单头IMU作同等强度的主张。
两种配置的通用jerk变化区间均跨零，因此walking subset上的平滑收益不能外推到全部动作。

作为进一步消融，Direct-V3在帧内增加sensor attention，并把batch标量ranking改为同传感器配置的逐样本paired/shuffled/zero ranking。
其512条预筛在validation上保持paired优于shuffled，但test-20仅有head 11/20、wrists 8/20改善，wrists平均误差反而增加0.0056 m，因此按预设门槛停止全量训练。
该负结果表明，简单保留传感器身份和单步ranking仍不足以解决生成过程中的实例级控制。

我们进一步尝试仅在训练期使用的Direct-V4 condition matcher。
它将IMU control token与目标motion投影到共享embedding，并在相同sensor config内采用双向InfoNCE直接辨识paired与其他真实实例。
512/128条train/validation预筛中，validation top-1从第一轮的56.4\%提升到第二轮的61.9\%，高于同配置batch内约25\%的随机水平；paired相对zero的margin满足率为96.9\%。
这表明实例信息可以被识别，但仍未达到预设70\%门槛，且matcher检索提升尚不能证明生成动作控制同步改善。
因此我们停止全量训练，并将其作为“辅助条件辨识并不自动转化为生成控制”的负结果。

需要注意的是，MDM Text-only 在某些 joint-proxy control 指标上可能并不比 ITM 差，因为它没有接收目标 IMU，误差只反映生成动作与某个目标 joint 轨迹的偶然接近程度；head/wrist joint proxy 也不能完全代表真实 IMU orientation consistency。
因此本文不把“所有 proxy error 都低于 Text-only”作为主张。
更可靠的结论来自受控条件交换：固定文本、随机噪声和长度，只改变 IMU 时，生成结果在轨迹、节奏和运动细节上发生可解释变化；同时 Stage-2b 相比 Stage-1 显著降低 jerk，使这种控制更平滑。

\begin{figure*}[t]
\centering
\fbox{\parbox[c][25mm][c]{0.94\textwidth}{\centering 图示占位：待插入最终实验可视化}}
\caption{固定文本\texttt{\protect\seqsplit{a person walks}}和扩散噪声，对比MDM Text-only、ITM Text+head IMU、ITM Text+wrists IMU，并在下方同步绘制目标IMU曲线和active-joint轨迹。候选来自\texttt{\protect\seqsplit{same\_text\_different\_imu/test\_walk\_head\_wrists}}。}
\label{fig:qualitative-1}
\end{figure*}

\subsubsection{Same Head IMU, Different Text}

固定 head-only IMU，使用 6 条 walking prompts。
Stage-1 的 prompt-level 结果如下：

\begin{table*}[t]
\centering
\caption{固定头部 IMU 时的文本条件对比}
\label{tab:result-14}
\footnotesize
\setlength{\tabcolsep}{3pt}
\renewcommand{\arraystretch}{1.15}
\begin{tabularx}{\linewidth}{LLLLL}
\toprule
Prompt & Head error & Arm swing & Jerk & Root travel \\
\midrule
walking with swinging arms & 0.1662 m & 0.2615 m & 6.2430 & 3.3377 m \\
walking quickly & 0.1824 m & 0.2304 m & 5.6510 & 2.4707 m \\
turning while walking & 0.2088 m & 0.2153 m & 4.1204 & 1.4661 m \\
walking & 0.1731 m & 0.2088 m & 5.1645 & 2.4170 m \\
walking slowly & 0.1776 m & 0.1742 m & 4.1510 & 2.1081 m \\
walking with large steps & 0.1948 m & 0.1582 m & 3.7824 & 1.8873 m \\
\bottomrule
\end{tabularx}
\end{table*}


Arm-swing prompt 的 arm swing proxy 最高，同时 head error 保持较低。
这比初始单样本 pilot 更支持文本补全未观测身体语义的假设。
不过，该 prompt 的 jerk 也最高，说明语义增强和动作平滑之间存在 trade-off。

Stage-2b在该实验上的结果更保守：active head error从0.1838 m升至0.2152 m，jerk ratio从4.8521升至5.6981，arm swing proxy从0.2080 m降至0.1581 m。
因此Stage-2b作为IMU-control和平滑主模型，而Stage-1作为head-only文本补全的诊断参考。
Stage-4 text-anchor将same-head active head error小幅改善到0.2103 m，arm swing proxy提升到0.1666 m，但jerk仍为5.9120，且arm swing相比Stage-1仍下降约19.9\%。
该结果构成本文讨论的主要trade-off。

为避免只依赖平均proxy和人工观察，我们使用修正后的变长采样器对30条test motions重新生成六种文本条件，每条动作保留自身39--196帧长度。
固定head IMU时，\texttt{\seqsplit{swinging arms}}相对普通walking的摆臂幅度增加$0.0979$ m（95\% CI $[0.0812,0.1144]$），29/30样本方向正确；\texttt{\seqsplit{turns while walking}}的累计root yaw增加$2.3689$ rad（95\% CI $[1.5691,3.1909]$），25/30方向正确。
进一步检查显示净yaw在26/30样本增加，净转角占累计转角的效率在23/30增加，说明该结果不只是高频yaw抖动。
相反，quick相对slow的路径速度差为$-0.0199$ m/s且区间跨零，仅5/30方向正确，cadence也仅19/30增加；large-steps的步幅proxy差为$-0.0237$ m且区间跨零，仅12/30方向正确，脚间距也仅14/30增加。
这与人工判断一致：摆臂与转向是当前较稳定的文本补全属性，快慢和大步尚未被可靠控制。
五类文本编辑使head error平均增加$0.0074$ m（95\% CI $[0.0031,0.0117]$），显示文本补全会轻微扰动局部IMU约束。

\begin{figure*}[t]
\centering
\fbox{\parbox[c][25mm][c]{0.94\textwidth}{\centering 图示占位：待插入最终实验可视化}}
\caption{固定head IMU和扩散噪声，比较\texttt{\protect\seqsplit{walks}}、\texttt{\protect\seqsplit{walks while swinging arms}}与\texttt{\protect\seqsplit{turns while walking}}。正文展示成功样例，补充材料同时保留快慢、大步等失败或弱响应样例。}
\label{fig:qualitative-2}
\end{figure*}

\subsubsection{Guidance Sweep}

扫描：

\[\begin{aligned}w_{\mathrm{text}}&\in\{1.0,2.5,4.0\},\\w_{\mathrm{imu}}&\in\{0.5,1.0,2.0\},\\\mathrm{sensor}&\in\{\mathrm{head},\mathrm{wrists}\}.\end{aligned}\]

用于定性展示的guidance配置为：

\begin{table*}[t]
\centering
\caption{Controlled Condition-Swapping Suite on Test Split：实验结果}
\label{tab:result-15}
\footnotesize
\setlength{\tabcolsep}{3pt}
\renewcommand{\arraystretch}{1.15}
\begin{tabularx}{\linewidth}{LLLLL}
\toprule
Sensor & Text scale & IMU scale & Active sensor error & Jerk \\
\midrule
head & 2.5 & 0.5 & 0.2003 m & 1.6134 \\
wrists & 2.5 & 0.5 & 0.4794 m & 1.5036 \\
wrists & 2.5 & 1.0 & 0.4847 m & 1.5066 \\
\bottomrule
\end{tabularx}
\end{table*}


我们进一步在共享的100-motion test subset上固定\texttt{\seqsplit{text\_scale=2.5}}、\texttt{\seqsplit{imu\_scale=1.0}}，扫描factorized guidance的\texttt{\seqsplit{joint\_scale}}为0、0.25、0.5、1和1.5。
Head各档active error为0.2131--0.2150 m，优于严格Text-only的样本数为58--61/100；wrists各档为0.4833--0.4868 m，改善样本数为69--71/100。
Wrists在\texttt{\seqsplit{joint\_scale=0}}时平均误差最低，但相对默认\texttt{\seqsplit{joint\_scale=1}}仅降低0.0022 m，95\%配对bootstrap区间为[-0.0058, 0.0015]，差异不显著；其余档位相对默认值的tracking与jerk区间也均跨零。
该结果说明当前控制主要来自IMU main effect，显式Text-IMU interaction scale在该范围内影响较弱，因此本文保留未针对测试集调优的默认值1.0。

\subsubsection{Matrix Experiment}

Matrix 实验随机选择 20 组 A/B test samples，对 head 和 wrists 分别生成 Text/IMU 的 8 种非空组合，并展示 Ground Truth A/B。
该实验主要用于定性分析：

\begin{itemize}
\item Text A + IMU B 是否表现为 A 的语义与 B 的运动细节。
\item Text B + IMU A 是否出现可解释的条件交换变化。
\item Text-only 是否保持 MDM 语义能力。
\item IMU-only diagnostic 是否提供轨迹/节奏但语义不足。
\end{itemize}

该矩阵直接呈现Text A + IMU B能否兼具A的语义与B的运动细节，也用于识别某一模态被模型忽略的失败情况。

\begin{figure*}[t]
\centering
\fbox{\parbox[c][25mm][c]{0.94\textwidth}{\centering 图示占位：待插入最终实验可视化}}
\caption{Ground Truth A/B位于首行，下方为Text None/A/B与IMU None/A/B的八种非空组合。候选使用Stage-2b的低jerk或高对比head/wrists矩阵样例。}
\label{fig:qualitative-3}
\end{figure*}

\subsection{Qualitative Analysis}

人工检查Figure 1--3候选后，可以稳定观察到文本和IMU都会改变生成结果，但控制强度并不均匀。
在固定head IMU的文本编辑中，\texttt{\seqsplit{walks while swinging arms}}和\texttt{\seqsplit{turns while walking}}通常比普通walking产生更清楚的上肢或方向变化；相反，\texttt{\seqsplit{walks slowly}}、\texttt{\seqsplit{walks quickly}}和\texttt{\seqsplit{walks with large steps}}在不同样本上的区分度不稳定。
部分动作能体现快慢或大步差异，但尚未形成跨样本一致现象。

同文本换IMU的结果更稳定：更换head或wrists信号会改变root travel、局部轨迹和身体摆动，同时Stage-2b较Stage-1显著减少高频晃动。
不过，现有样例主要证明“条件确实产生影响”，而不是每个条件都能实现精确语义或传感器匹配。
因此正文图应同时呈现成功与失败案例，避免仅凭挑选样例扩大结论。

MotionLab \cite{guo2025motionlab} test-20中按active trajectory proxy选出的best/worst样例也体现同一问题：部分metric-best样例降低了局部轨迹误差，却削弱了“抬手走路”“跌倒后起身”等文本语义；跳跃开合和肩高拍臂只得到部分形态。
抽帧检查未发现持续骨长变化或早期回归模型式高频果冻抖动，但没有样例同时达到文本语义清楚、IMU影响可解释和视觉自然三个要求。
因此MotionLab图只作为迁移诊断，不以metric-best替代人工语义判断。

\subsection{Preliminary Zero-shot Study on Real Nymeria IMUs}

为检查从AMASS/SMPL虚拟IMU训练得到的控制器能否接收真实传感器，我们在Nymeria上进行零样本预备实验。
Nymeria narration的\texttt{\seqsplit{start\_time/end\_time}}属于head Aria的\texttt{\seqsplit{DEVICE\_TIME}}秒，Xsens \texttt{\seqsplit{timestamps\_us}}属于公共\texttt{\seqsplit{TIME\_CODE}}微秒。
我们使用head VRS提供的设备时间到timecode转换，将每条narration精确索引到\texttt{\seqsplit{body/xdata.npz}}，再把Xsens Z-up坐标转换到ITM的Y-up六槽表示。
实验选取3条来自不同录制序列、约5秒且head/wrists四元数有效的片段；另外5条候选因目标传感器四元数无效而排除。

模型使用Stage-2b且不进行Nymeria微调，MDM Text-only、ITM+head与ITM+wrists共享caption、长度、seed和初始噪声：

\begin{table*}[t]
\centering
\caption{Nymeria 真实 IMU 零样本预备实验}
\label{tab:result-16}
\footnotesize
\setlength{\tabcolsep}{3pt}
\renewcommand{\arraystretch}{1.15}
\begin{tabularx}{\linewidth}{LLLLLL}
\toprule
Input & Active sensor trajectory error \ensuremath{\downarrow} & Root-relative error \ensuremath{\downarrow} & Jerk ratio \ensuremath{\downarrow} & Arm swing & Root travel \\
\midrule
MDM Text-only & 0.365 m & 0.294 m & 9.065 & 0.240 m & 0.402 m \\
ITM + head & \textbf{0.198 m} & \textbf{0.241 m} & \textbf{7.789} & 0.147 m & 0.290 m \\
ITM + wrists & 0.462 m & 0.344 m & 14.247 & 0.320 m & 0.599 m \\
\bottomrule
\end{tabularx}
\end{table*}


Head配置的active trajectory error相对Text-only降低约45.8\%，说明真实head信号在零样本设置下没有被完全忽略。
Wrists配置则同时恶化轨迹误差和jerk，反映真实Xsens与训练时SMPL虚拟腕部IMU之间存在安装朝向、噪声和统计分布偏移。
由于样本仅3条，且生成骨架与Nymeria Xsens骨架并不完全一致，该结果只作为可行性诊断，不进入主表，也不能视为对Ego4o或其他Nymeria方法的正式比较。

\begin{figure*}[t]
\centering
\fbox{\parbox[c][25mm][c]{0.94\textwidth}{\centering 图示占位：待插入最终实验可视化}}
\caption{每条Nymeria片段展示GT、MDM Text-only、ITM+head和ITM+wrists，以及同步真实IMU曲线；现有可视化位于\texttt{\protect\seqsplit{outputs/nymeria\_pilot/stage2b\_zero\_shot/per\_clip/}}。}
\label{fig:qualitative-4}
\end{figure*}

% !TeX root = ../main-cn.tex
% Converted from paper/论文草稿.md; UTF-8.
\section{讨论与限制}
\label{sec:5_discussion}


实验支持将文本和IMU视为两类互补但相互竞争的条件。
固定文本和扩散噪声时，更换IMU能够改变生成动作的轨迹、节奏和摆动，说明控制器并非仅依赖文本先验采样。
Stage-2b进一步表明，无需解冻MDM即可通过轻量trajectory、velocity和jerk目标改善时间稳定性。
与此同时，5-rep子集评估显示head-only ITM基本保持文本检索能力，但FID有所上升，表明控制收益伴随生成分布偏移。
MotionLab迁移实验进一步显示，更强生成骨干能显著改善文本生成指标，但不会自动解决IMU实例级控制；平均误差改善与逐样本稳定可控是两个不同要求。

更值得关注的是IMU控制与文本补全之间的trade-off。
Stage-1在head-only条件下呈现较强摆臂响应，却具有较高jerk；Stage-2b显著平滑同文本换IMU结果，但其same-head arm-swing proxy下降；Stage-4利用frozen MDM Text-only prediction进行anchor，仅恢复部分上肢变化。
这说明简单joint-space约束难以同时保证局部传感器一致性、全身自然度和未观测部位语义。
后续方法可能需要保留传感器身份的跨部位交互、分层语义控制或更强生成骨干，而不只是继续调整三个辅助损失权重。

本文存在四项主要限制。
第一，active-joint trajectory与acceleration仍是joint-space proxy。
我们在5条GT动作上以200次优化拟合neutral SMPL：head/wrists trajectory误差仅为0.00765/0.00781 m，但virtual acceleration误差仍为0.885/2.564 m/s2，orientation geodesic error为0.320/1.032 rad，均未通过预设门槛。
该结果说明HumanML 22 joints无法可靠恢复骨轴twist，尤其是腕部方向，因此本文不把拟合后的generated IMU用于正式主张。
第二，generation-quality结果来自682条可用动作中的前672条和5次重复，不能替代full-set 20-rep HumanML3D协议。
第三，定性实验中摆臂和转向比快慢、大步等细粒度语义稳定，模型尚不能精确解耦所有文本属性。
第四，Nymeria仅包含3条零样本片段，wrists存在明显域偏移。
正式真实IMU实验需要subject-disjoint划分、每传感器静态安装校准、真实/虚拟IMU统计匹配，以及仅微调IMU encoder/adapters的对照。

外部比较也受任务协议限制。
MDM是同一生成骨干下的直接Text-only baseline，Flexible IMUPoser提供IMU-only reconstruction边界；Ego4o与Spatial-Related Sensors Matters需要在统一Nymeria划分和相同输入配置下重新训练，才能使用本文的生成与控制指标公平比较。
本文因此将结论限定于任务可行性、轻量融合方法和可观测trade-off，而不声称超过所有Text-to-Motion或IMU姿态估计方法。

% !TeX root = ../main-cn.tex
% Converted from paper/论文草稿.md; UTF-8.
\section{结论}
\label{sec:6_conclusion}


本文提出IMU-guided Text-to-Motion Generation，将稀疏IMU作为预训练文本动作生成器的实例级控制条件。
所提出的方法冻结MDM和CLIP，仅通过时序IMU encoder与逐层零初始化residual adapters学习控制，在较低训练成本下保留原有文本生成路径。
HumanML3D/AMASS实验表明，IMU可以在模糊文本下改变具体运动细节，正则化主模型能够显著降低该过程中的jerk，同时基本保持文本检索指标。
实验也揭示了局部IMU一致性、动作平滑性和未观测身体语义补全之间的系统性权衡；简单text anchor尚不足以同时优化三者。
Nymeria真实IMU预备实验显示head控制具有初步跨域可行性，但wrists域偏移仍需校准。
上述结果为结合语言语义与可穿戴传感器的可控人体动作生成提供了任务定义、轻量实现和受控条件交换评价基础。

\bibliographystyle{IEEEtran}
\bibliography{sec-cn/references}
\end{document}
